\section{Fault Causality Analysis}
\label{sec:fia}



This section details 
the types of faults to inject (\Cref{sec:fia:faults-to-inject}), 
the fault injection process (\Cref{sec:fia:how-to-inject}), 
and the approach for establishing causal relationships between faults (\Cref{sec:fia:establish-causality}). 

\subsection{What faults to inject?} 
\label{sec:fia:faults-to-inject}

A recent study~\cite{qian2023vicious} 
on \vcs from open-source systems 
reveals that \vcs are often caused by 
functional and performance interferences 
between requests and error handlers. 
Functional interferences typically manifest 
as exceptions and errors captured by system-specific error detectors, 
and performance interferences typically manifest as contention. 
Following this observation, 
\tool injects three types of faults. 

\textbf{Exception}. 
A target system can encounter three types of exceptions 
during its execution: 
(1) system-specific exceptions 
-- exceptions thrown explicitly inside the target system's source code, 
(2) library function exceptions 
-- exceptions thrown by a library or native function, 
(3) unchecked (runtime) exceptions 
-- exceptions that are not required to be explicitly handled 
(by try/catch blocks in JVM-based systems) 
or declared explicitly in function signatures. 

\tool injects system-specific exceptions 
and library function exceptions, 
because they occur at limited, explicitly-declared program locations, 
while unchecked exceptions can happen at arbitrary locations. 
Note that \tool still injects an unchecked exception 
if it is explicitly thrown in the system's code 
-- e.g., an \texttt{IllegalArgumentException} is often thrown 
when a request with invalid input is received. 
In practice, 
\tool filters out reflection-related and
security-related exceptions, 
as they tend to terminate or hang the system 
rather than propagate. 
\tool also ignores exceptions 
only reachable from tests. 


\textbf{Contention}. 
Though contention can happen at arbitrary program locations, 
\tool only considers loops for contention injection 
to simulate resource-intensive tasks causing contention, 
due to their association with performance issues discovered by existing studies~\cite{li2018pcatch, jin2012understanding, qian2023vicious}.
\tool uses iteration count of the workload-related loops to capture contention induced by increased workload (\Cref{sec:fia:establish-causality}), similar to the approach adopted by \citet{li2018pcatch}. 

\tool adopts a loop scalability analysis 
to filter out loops unlikely to cause performance issues. 
\tool excludes loops with a constant upper bound on their iteration count, detected by a best-effort data flow analysis to track the loop guard condition to a constant. 
In addition, 
\tool ranks the size of code reachable recursively starting from a loop 
and identifies loops involving I/O operations, 
to exclude loops that have a short execution (i.e., lowest ranked 10\% of loops) and do not perform I/O. 

\textbf{System-Specific Error}. 
Many distributed systems implement system-specific error detectors, 
such as status checks and health monitoring, 
without using exceptions. 
Such system-specific error detectors 
are often implemented as functions with a boolean return value. 
For example, in HDFS, a NameNode thread checks 
if a DataNode's heartbeat is received within a timeout limit 
using a function \texttt{node.isStale()}.

\tool filters out boolean-returning functions 
if they are unlikely to be a system-specific error detector. 
For example, 
\tool filters out such boolean-returning functions in JDK libraries, 
such as \texttt{contains()} function of Java collection type objects 
(detailed filtering criteria in \Cref{sec:impl-detail}). 
\tool's fault filtering criteria 
is designed to be conservative
to avoid meaningful fault being skipped accidentally. 
The injection points that do not impact system execution will be automatically deprioritized during fault injection experiments 
by the test budget allocation protocol (\Cref{sec:alloc}). 


\subsection{How to inject faults?}
\label{sec:fia:how-to-inject}

\begin{figure}[tb]
\centering
\footnotesize
\begin{minted}[
  xleftmargin=2ex, xrightmargin=1ex, frame=lines, fontsize=\footnotesize, linenos, 
  autogobble, numbersep=1ex, framesep=1mm, breaklines]
{java}
for (int i = 0; i < 10; i++) { // Loop 1
  sock.read(); // TP1
  if (condition1) {// TP2; MP1
    throw new IOException();
  }
  for (int j = 0; j < 10; j++) { // Loop 2
    if (foo()) { // TP3; MP2
      throw new IOException();
    }
  }
}

boolean foo() { // NP1
  fis = new FileInputStream(path); // TP 4
  if (fis.read() == -1) { // MP 3; TP 5
    return false;
  }
  return true;
}
\end{minted}
% \vspace{-3mm}
\caption{Pseudo-code demonstrating the injection and monitor points. ``TP'' means throw point and ``NP'' negation point. 
``MP'' means monitor point (used and explained in \Cref{sec:chain-forming:approx-check}).}
% \vspace{-1mm}
\label{fig:injection-point-exp}
\end{figure}


\noindent
\tool's static analyzer 
identifies locations to inject faults 
and instruments fault injection hooks statically. 
The hook transfers the control to \tool's runtime agent. 
During fault injection experiments, 
\tool dynamically injects faults 
when corresponding hooks are encountered.  
We use 
a Java-style pseudo-code with eight fault locations
(shown in \Cref{fig:injection-point-exp})
to explain how fault injection is performed. 




\textbf{Exception Injection}. 
For system-specific exceptions, 
as they are typically thrown inside an if-statement 
(lines 4 and 8), 
\tool injects 
a one-time throw of the same exception 
when the 
if-statement (lines 3 and 7, referred to as \textit{Throw Point}) is reached. 
For library function exceptions, 
\tool injects a one-time throw of the exception 
declared in the function signature 
at the invocation site of these functions 
(lines 2, 14, and 15). 
Exceptions are constructed at runtime 
using their simplest constructor.

\textbf{Contention Injection}. 
For potentially non-scalable loops 
identified by \tool (e.g., line 1 and 6), 
\tool 
injects a predefined amount of \textbf{spinning delay} 
into each loop iteration 
before the first line of the loop body executes 
to simulate potential contention induced 
by executing a large number of iterations of this loop. 
Due to the way we inject contention, 
we refer to such injection as \textit{delay injection}. 


Each delay injection is performed seven times with varying length (100ms to 8s) to maximize the discovery of causal relationships between faults. 
The system is configured with \textit{reduced timeouts} (10-20 seconds) to make it more sensitive to additional workload.


The delay and timeout values are determined empirically to maximize the impact of delay injections. We first lower system timeout configurations to a safe range of 10-20s via trial and error, which ensures normal system functionality being preserved (i.e., all unit and integration tests still pass).
Based on this threshold, we select seven static delay values between 100ms and 8s that are likely to cause timeouts when injected repeatedly inside loops. 
We identify relevant configuration items using the keywords of ``timeout'' and ``interval''. 
The timeout adjustment process typically takes a student 30 minutes per system, and are not essential for \tool's effectiveness, as delay-related causal relationships can still be observed with default settings.


\textbf{System-Specific Error Injection}. 
Because system-specific errors 
are captured by system-specific error detectors, 
whose return values are boolean typed, 
\tool negates the return value of these functions 
(which is referred to as \textit{Negation Point}) 
before the execution returns to its caller 
to simulate the effect of such faults 
(\texttt{foo()} at line 13). 
We refer to the injection of system-specific error as a \textit{negation injection}. 



\subsection{How to establish causal relationships?} 
\label{sec:fia:establish-causality}

\label{sec:fia:interference-analysis}

\tool records encountered faults during the injection run 
and profile run, 
and compares them to identify additional faults triggered 
to establish causal relationships. 
We categorize the encountered faults into the following categories: 

\begin{enumerate}[wide, labelwidth=!, labelindent=\parindent, noitemsep, topsep=0pt]
\item \textbf{Execution Trace Interference}. 
For \textbf{exceptions}, 
\tool monitors whether the throw statement is reached. 
For \textbf{system-specific errors}, 
\tool monitors whether 
the return value of the error-detector function is negated. 

\item \textbf{Iteration Count Interference}. 
For \textbf{contention}, 
\tool monitors whether any loop's iteration count \textit{statistically} increases compared to the profile run, 
as an indicator for contention. 
The loop iteration count serves as a good indicator for the amount of workload processed~\cite{li2018pcatch}. 
We use one-sided t-test~\cite{student1908probable} with a $p$-value of 0.1 for statistical significance.

\end{enumerate}

\tool executes each profile run and injection run five times to reduce the impact of \textbf{non-determinism} in the system execution. This also allows us to use statistical tests on loop iteration counts.


Injecting $f_1$ may trigger multiple additional faults, stemming from both interference types. After all the injection runs for $f_1$, \tool collects a list of additional faults triggered $[f_2, f_3, ..., f_n]$. 



\textbf{Nested and Consecutive Loops.}
One special case of iteration count interference is the handling of nested and consecutive loops in the system. 
Three types of compositions of workload related loops exist: 1) independent loops, 2) nested loops, and 3) consecutive loops. 
We observe that 2) nested and 3) consecutive workload-related loops are often used in batch processing. 
In such loops, 
% where 
a delayed sub-request 
can propagate delays to its parent request (the batch) 
and the next sub-request in the same batch.  
For example, in a batched RPC request, a delayed sub-request can time out the entire call as well as the next sub-request.

\begin{figure}
\centering
\footnotesize
\begin{minted}[
  xleftmargin=2ex, xrightmargin=1ex, frame=lines, fontsize=\footnotesize, linenos, 
  autogobble, numbersep=1ex, framesep=1mm, breaklines]
{java}
while (shouldRun()) { // Loop 1
    resp = sendHeartBeat();
    // Loop 2
    for (DataNodeCommand cmd: resp.getCommands()) {
        processCommand(cmd);
    }
    // ...
    List<BlockReport> reports = new ArrayList<>();
    // Loop 3
    for (Map.Entry kv: perVolumeBlocks.entrySet()) { 
        reports.add(convertFormat(kv));
    }
    // ...
    for (BlockReport report: reports) {
        nnRpcStub.blockReport(report);
        // ...
    }
}
\end{minted}
% \vspace{-3mm}
\caption{Pseudo-code demonstrating delay in nested loops, code simplified from \texttt{BPServiceActor} in HDFS.}
\label{fig:nested-loop-delay}
\end{figure}


\tool identifies the parent-child loop pairs 
to establish this causal relationship and expand the impact scope of contention injection.
\Cref{fig:nested-loop-delay} shows the pseudo-code demonstrating contentions in the nested loops. 
If injection $f_1$ increases iterations in loop 2 ($L_2$), its parent ($L_1$) and sibling ($L_3$) can also be affected.
We use \CFGBackEdge{$L_2$}{$L_1$} to represent the former (``I'' for ``Inverse'') and \CFGEdge{$L_1$}{$L_3$} to represent the latter. 

\textbf{Summary for Fault Causality}. 
\Cref{tab:fault-interference-summary} summarizes all the six causal relationships between different types of faults. 
The first four are differentiated 
by the injected fault and the additional fault triggered. 
The last two are special causal relationships handling nested loops, extending the impact range of contentions. 


\begin{enumerate}[leftmargin=3ex, noitemsep, topsep=2pt]
\item \DelayExecEdge{$f_1$}{$f_2$}: Injecting a delay $f_1$ into a loop causes an additional exception or negation $f_2$. Delay $f_1$ has an execution trace interference $f_2$ on the system (hence the name \texttt{E(D)}).

\item \DelayIterEdge{$f_1$}{$f_2$}: Injecting a delay $f_1$ into a loop causes an additional delay ($f_2$) in another loop. Delay $f_1$ has an iteration count interference $f_2$ on the system (name \texttt{S$^{+}$} indicates a statistically significant increase).

\item \FaultExecEdge{$f_1$}{$f_2$}: Injecting an exception or negation $f_1$ causes an additional exception or negation $f_2$. 

\item \FaultIterEdge{$f_1$}{$f_2$}: Injecting an exception or negation $f_1$ causes an additional delay ($f_2$) in another loop.

\item \CFGBackEdge{$f_2$}{$f_2'$}: $f_2$ is an additional delay caused by the injection $f_1$, affecting its parent loop ({$f_2'$}).

\item \CFGEdge{$f_2'$}{$f_2''$}: Following \CFGBackEdge{$f_2$}{$f_2'$}, the parent loop delay ($f_2'$) further affects the sibling loop ($f_2''$).
\end{enumerate}

\begin{table}[t]
\centering
\footnotesize

\caption{Summary of the causal relationships between faults.}
% \vspace{-3mm}

\begin{tabular}{ccc}
\toprule
\textbf{Type} & \textbf{Injected Fault} & \textbf{Additional Fault} \\
\midrule

\DelayExecEdge{}{} & Delay & Exception; Negation \\
\DelayIterEdge{}{} & Delay & Delay \\
\FaultExecEdge{}{} & Exception; Negation & Exception; Negation \\
\FaultIterEdge{}{} & Exception; Negation & Delay \\

\midrule 
\CFGBackEdge{}{} & \multicolumn{2}{c}{$f_1\longSPlusArrow{D/I}$\CFGBackEdge{$f_2$}{$f_2'$} only} \\
\CFGEdge{}{} & \multicolumn{2}{c}{$f_1\longSPlusArrow{D/I}$\CFGBackEdge{$f_2$}{$f_2'$}\CFGEdge{}{$f_2''$} only} \\

\bottomrule

\end{tabular}
\vspace{-2mm}

\label{tab:fault-interference-summary}
\end{table}

